\begin{pastexam}[label={px:gw3-110-ds-04}]{關務三等資訊處理 110 年 第四題}
（申論）

密文（Cipher text or Cypher text）：\underline{明}請到家玩\underline{天}你我來，應用環狀串列（circular linked list），請問明文
（Plain text or Clear text）為何？（15 分）
\end{pastexam}

\begin{pastsoln}
答案：明文是「明天請你到我家來玩」。

\textbf{題目的解讀。} 題目只給密文和「環狀串列」，沒有說明規則；畫底線的「明」「天」是提示：明文的前兩個字是
「明天」。 把 $9$ 個字依序放進環狀串列，「明」是第 $0$ 個節點，「天」是第 $5$ 個節點，從「明」走 $5$ 步到「天」。
所以解讀為：\textbf{從「明」開始，每次沿串列往後走 $5$ 個節點，讀出經過停下的字}，讀滿 $9$ 個字為止。

\textbf{環狀串列。} 最後一個節點的 \textit{next} 指回第一個節點（資結書第~5.2~節），走到尾端會自動繞回開頭，正好
用來做「往後數 $5$ 個、超過尾端就從頭繼續」：
\begin{center}
\begin{tikzpicture}[nd/.style={circle,draw,minimum size=8mm,inner sep=0pt,font=\small},
  ar/.style={-{Stealth[length=1.6mm]}}]
\foreach \c [count=\i from 0] in {明,請,到,家,玩,天,你,我,來} {
  \node[nd] (n\i) at ({90 - 40*\i}:2.1) {\c};
  \node[font=\scriptsize] at ({90 - 40*\i}:2.85) {\i};
}
\foreach \i [evaluate=\i as \j using {int(mod(\i+1,9))}] in {0,...,8} {
  \draw[ar] (n\i) to[bend left=12] (n\j);
}
\end{tikzpicture}
\end{center}

\textbf{解密過程。} 目前位置 $p$ 從 $0$ 開始，每讀一個字就 $p \gets (p + 5) \bmod 9$（在環狀串列上就是沿 \textit{next}
走 $5$ 次）：
\begin{center}
\begin{tabular}{lccccccccc}
\toprule
第幾個字 & 1 & 2 & 3 & 4 & 5 & 6 & 7 & 8 & 9 \\
節點位置 $p$ & 0 & 5 & 1 & 6 & 2 & 7 & 3 & 8 & 4 \\
讀出 & 明 & 天 & 請 & 你 & 到 & 我 & 家 & 來 & 玩 \\
\bottomrule
\end{tabular}
\end{center}
因為 $\gcd(5, 9) = 1$，走 $9$ 次恰好把每個節點各經過一次，不會重複也不會漏掉，得到「明天請你到我家來玩」，是通順
的句子，印證了上面的解讀。

\textbf{加密是怎麼做的。} 反過來看，密文是把明文的第 $1, 3, 5, 7, 9$ 個字（明、請、到、家、玩）寫在前面，再接第
$2, 4, 6, 8$ 個字（天、你、我、來），也就是兩列的「柵欄密碼」（rail fence）。 在長度 $9$ 的環上，明文相鄰兩字在
密文中相距 $5$ 個位置，所以解密就是每次跳 $5$ 格。 注意這裡節點讀過之後\textbf{不刪除}；若改成像 Josephus 問題那樣
讀一個刪一個再數 $5$ 個，第三個字會讀到「到」，句子不通。
\end{pastsoln}
